\chapter{QUẢN LÝ TOPOLOGY}
\label{chapter:chapter-5}
% Yêu cầu package: \usepackage{tikz} và \usetikzlibrary{arrows.meta,shapes,positioning}

% ============================================================
\section{Kiến Trúc Và Khám Phá Topology}
% ============================================================

Mạng Multi-AP được tổ chức theo \textbf{kiến trúc hình cây} với một Controller duy nhất
ở gốc và các Agent tại các nút lá. Để xây dựng cấu trúc này, hệ thống sử dụng các bản
tin mở rộng của IEEE~1905.1 nhằm giải quyết hai bài toán cốt lõi: (1)~tìm kiếm và xác
định vai trò giữa các thiết bị, và (2)~thu thập thông tin cấu trúc mạng.

Mỗi Agent khi khởi động sẽ chủ động phát quảng bá để tìm Controller, sau đó Controller
sẽ truy vấn ngược lại để lấy thông tin chi tiết về các BSS và các thiết bị khách đang
kết nối với Agent đó.


% ------------------------------------------------------------
\subsection{Luồng Bản Tin Khám Phá (Discovery Flow)}
% ------------------------------------------------------------

Cặp bản tin \textit{Search/Response} cho phép Agent tìm ra Controller duy nhất trong
mạng. Phiên bản \textit{Extended} bổ sung thêm các TLV mang thông tin về Profile bảo
mật (Profile-1, -2 hay -3) và loại thiết bị, giúp hai bên thỏa thuận cơ chế xác thực
sẽ sử dụng trước khi trao đổi cấu hình.

\begin{center}
	\begin{tikzpicture}[>=stealth, font=\small]
		% --- Tiêu đề ---
		\node[draw, fill=gray!15, minimum width=2.8cm, minimum height=0.6cm]
		(A) at (0, 0) {Multi-AP Agent};
		\node[draw, fill=gray!15, minimum width=2.8cm, minimum height=0.6cm]
		(C) at (8, 0) {Multi-AP Controller};
		
		% --- Lifelines ---
		\draw[dashed, gray!50] (0,-0.4) -- (0,-10.5);
		\draw[dashed, gray!50] (8,-0.4) -- (8,-10.5);
		
		% --- Bước 1 & 2: Search / Response ---
		\draw[->] (0,-1.2) -- (8,-1.2)
		node[midway, above, scale=0.7, fill=white, inner sep=1pt]
		{(1) 1905 AP-Autoconfig Search (Extended)};
		\node[scale=0.6, gray, align=center] at (4,-1.9)
		{\textit{Agent phát quảng bá --- mang thông tin Profile và loại thiết bị}};
		
		\draw[->] (8,-3.0) -- (0,-3.0)
		node[midway, above, scale=0.7, fill=white, inner sep=1pt]
		{(2) 1905 AP-Autoconfig Response (Extended)};
		\node[scale=0.6, gray, align=center] at (4,-3.7)
		{\textit{Controller xác nhận vai trò và cơ chế bảo mật sẽ áp dụng}};
		
		% --- Phân tách pha ---
		\draw[dashed, gray!60] (-1.5,-4.4) -- (9.5,-4.4);
		\node[scale=0.65, gray!80] at (4,-4.75)
		{Sau khi xác định vai trò --- Controller truy vấn chi tiết cấu trúc mạng};
		
		% --- Bước 3 & 4: Topology Query / Response ---
		\draw[->, blue] (8,-5.6) -- (0,-5.6)
		node[midway, above, scale=0.7, fill=white, inner sep=1pt]
		{(3) 1905 Topology Query [Profile TLV]};
		\node[scale=0.6, gray, align=center] at (4,-6.3)
		{\textit{Controller yêu cầu Agent báo cáo toàn bộ thông tin cấu trúc}};
		
		\draw[->, blue] (0,-7.4) -- (8,-7.4)
		node[midway, above, scale=0.7, fill=white, inner sep=1pt]
		{(4) 1905 Topology Response [BSS + Client Info]};
		\node[scale=0.6, gray, align=center] at (4,-8.1)
		{\textit{Agent trả về danh sách BSSID, kênh radio và thiết bị khách đang kết nối}};
	\end{tikzpicture}
\end{center}

\textbf{Giải thích luồng bản tin:}
\begin{enumerate}
	\item Agent phát quảng bá \textit{AP-Autoconfig Search} phiên bản Extended trên
	mạng 1905.1. Bản tin mang thêm TLV khai báo Profile bảo mật hỗ trợ
	(Profile-1/2/3) và loại thiết bị (Gateway, Repeater).
	\item Controller phản hồi \textit{AP-Autoconfig Response}, xác nhận vai trò điều
	phối và thông báo Profile cao nhất được chấp thuận cho phiên làm việc này.
	\item Controller chủ động gửi \textit{Topology Query} kèm \textit{Profile TLV} để
	yêu cầu Agent cung cấp toàn bộ thông tin cấu trúc hiện tại.
	\item Agent phản hồi bằng \textit{Topology Response} liệt kê tất cả BSS đang hoạt
	động (BSSID, kênh, băng thông) và danh sách thiết bị khách đang liên kết ---
	dữ liệu nền tảng để Controller xây dựng bản đồ mạng.
\end{enumerate}


% ============================================================
\section{Quản Lý Trạng Thái Thiết Bị Khách (Client Management)}
% ============================================================

Controller cần nắm được trạng thái kết nối của mọi thiết bị khách (STA) trong mạng
theo thời gian thực để đưa ra quyết định chuyển vùng, phân bổ tài nguyên và cập nhật
bảng định tuyến. Để đáp ứng yêu cầu này, hệ thống sử dụng cơ chế
\textbf{Fast Notification} --- Agent báo cáo ngay lập tức mỗi khi có sự thay đổi trạng
thái kết nối, thay vì chờ chu kỳ báo cáo định kỳ.


% ------------------------------------------------------------
\subsection{Luồng Bản Tin Sự Kiện Kết Nối Client}
% ------------------------------------------------------------

Mỗi khi một thiết bị khách (STA) gia nhập hoặc rời khỏi một điểm truy cập, Agent quản
lý BSS đó ngay lập tức phát bản tin thông báo đến Controller. Controller sau đó phát
quảng bá thông tin cập nhật đến tất cả các Agent còn lại để đồng bộ bảng định tuyến
toàn mạng.

\begin{center}
	\begin{tikzpicture}[>=stealth, font=\small]
		% --- Tiêu đề ---
		\node[draw, fill=gray!15, minimum width=2.2cm, minimum height=0.6cm]
		(STA) at (0, 0)   {Client (STA)};
		\node[draw, fill=gray!15, minimum width=2.2cm, minimum height=0.6cm]
		(A)   at (5.5, 0) {Agent};
		\node[draw, fill=gray!15, minimum width=2.8cm, minimum height=0.6cm]
		(C)   at (11, 0)  {Controller};
		
		% --- Lifelines ---
		\draw[dashed, gray!50] (0,   -0.4) -- (0,   -12);
		\draw[dashed, gray!50] (5.5, -0.4) -- (5.5, -12);
		\draw[dashed, gray!50] (11,  -0.4) -- (11,  -12);
		
		% --- Nhãn giai đoạn: Association ---
		\node[scale=0.72, font=\bfseries, anchor=west] at (-0.5,-1.0)
		{Sự kiện: STA gia nhập mạng};
		\draw[gray!60] (-0.5,-1.3) -- (12,-1.3);
		
		% --- Bước 1: STA Associate ---
		\draw[->, dashed] (0,-2.0) -- (5.5,-2.0)
		node[midway, above, scale=0.7, fill=white, inner sep=1pt]
		{(1) 802.11 Association};
		\node[scale=0.6, gray, align=left, anchor=west] at (0.2,-2.7)
		{\textit{STA hoàn tất thủ tục 802.11 Association / 4-way Handshake với Agent}};
		
		% --- Bước 2: Topology Notification (Association) ---
		\draw[->, red] (5.5,-3.6) -- (11,-3.6)
		node[midway, above, scale=0.7, fill=white, inner sep=1pt]
		{(2) 1905 Topology Notification [Client Join]};
		
		% --- Bước 3: Controller phát quảng bá ---
		\draw[->, red, dashed] (11,-4.6) -- (5.5,-4.6)
		node[midway, above, scale=0.7, fill=white, inner sep=1pt]
		{(3) Broadcast Update \textrightarrow{} All Agents};
		\node[scale=0.6, gray, align=center] at (8.25,-5.3)
		{\textit{Controller cập nhật bảng định tuyến và thông báo cho toàn mạng}};
		
		% --- Nhãn giai đoạn: Disassociation ---
		\node[scale=0.72, font=\bfseries, anchor=west] at (-0.5,-6.4)
		{Sự kiện: STA rời mạng};
		\draw[gray!60] (-0.5,-6.7) -- (12,-6.7);
		
		% --- Bước 4: STA Disassociate ---
		\draw[->, dashed] (0,-7.4) -- (5.5,-7.4)
		node[midway, above, scale=0.7, fill=white, inner sep=1pt]
		{(4) 802.11 Disassociation};
		
		% --- Bước 5: Topology Notification (Disassociation) ---
		\draw[->, brown] (5.5,-8.4) -- (11,-8.4)
		node[midway, above, scale=0.7, fill=white, inner sep=1pt]
		{(5) 1905 Topology Notification [Client Leave]};
		
		% --- Bước 6: Client Disassociation Stats ---
		\draw[->, brown] (5.5,-9.6) -- (11,-9.6)
		node[midway, above, scale=0.7, fill=white, inner sep=1pt]
		{(6) Client Disassociation Stats TLV};
		\node[scale=0.6, gray, align=center] at (8.25,-10.3)
		{\textit{Lý do ngắt kết nối, thời lượng phiên, lưu lượng đã truyền}};
		
		% --- Bước 7: Controller cập nhật ---
		\draw[->, brown, dashed] (11,-11.2) -- (5.5,-11.2)
		node[midway, above, scale=0.7, fill=white, inner sep=1pt]
		{(7) Broadcast Update \textrightarrow{} All Agents};
	\end{tikzpicture}
\end{center}

\textbf{Giải thích luồng bản tin:}
\begin{enumerate}
	\item STA hoàn tất thủ tục 802.11 Association và 4-way Handshake với Agent, thiết
	lập kết nối Wi-Fi ở lớp vật lý.
	\item Agent lập tức gửi \textit{1905 Topology Notification} đến Controller, thông
	báo có STA mới gia nhập kèm địa chỉ MAC và BSSID liên quan.
	\item Controller cập nhật bảng định tuyến nội bộ và phát quảng bá thông tin này
	đến tất cả các Agent còn lại trong mạng để đồng bộ.
	\item STA gửi khung \textit{Disassociation} hoặc bị Agent ngắt kết nối do timeout.
	\item Agent ngay lập tức gửi \textit{1905 Topology Notification} thông báo STA đã
	rời mạng, kèm lý do ngắt kết nối.
	\item Agent gửi thêm \textit{Client Disassociation Stats TLV} cung cấp số liệu
	phiên kết nối: thời lượng, tổng lưu lượng truyền nhận và lý do ngắt kết nối
	chi tiết --- dữ liệu này giúp Controller tối ưu hóa chính sách chuyển vùng
	cho các lần kết nối tiếp theo.
	\item Controller phát quảng bá cập nhật đến tất cả Agent, đảm bảo bảng định tuyến
	toàn mạng được đồng bộ.
\end{enumerate}


% ============================================================
\section{Báo Cáo Topology Cho Wi-Fi 7 (MLO)}
% ============================================================

Wi-Fi 7 (IEEE 802.11be) giới thiệu kiến trúc \textbf{Multi-Link Operation (MLO)}, cho
phép một thiết bị duy trì đồng thời nhiều liên kết trên các băng tần khác nhau. Điều
này làm phức tạp đáng kể cấu trúc topology so với các thế hệ trước: thay vì ánh xạ
đơn giản \textit{một STA $\leftrightarrow$ một BSSID}, hệ thống nay phải quản lý các
thực thể \textit{Multi-Link Device (MLD)} bao gồm nhiều địa chỉ MAC và liên kết
đồng thời.

\begin{itemize}
	\item \textbf{MLD Structure} (\textit{Cấu trúc AP MLD}): Agent báo cáo địa chỉ
	MAC của thực thể AP MLD cùng với danh sách các \textit{Affiliated AP} --- các
	radio thành viên thuộc cùng một AP MLD. Mỗi Affiliated AP có địa chỉ BSSID
	riêng nhưng chia sẻ chung một địa chỉ MLD MAC ở lớp quản lý. Controller cần
	thông tin này để hiểu radio nào thuộc cùng một thực thể logic và điều phối
	phân bổ lưu lượng giữa chúng.
	
	\item \textbf{Associated STA MLD} (\textit{STA đa liên kết đang kết nối}): Agent
	báo cáo chi tiết các \textit{Setup Link} --- tập hợp các liên kết được thiết
	lập giữa một Client MLD và AP MLD cụ thể. Mỗi Setup Link tương ứng với một
	cặp (link ID, BSSID, địa chỉ MAC STA trên băng tần đó). Thông tin này giúp
	Controller biết chính xác thiết bị nào đang dùng bao nhiêu liên kết đồng thời
	và trên băng tần nào, từ đó đưa ra quyết định chuyển vùng hoặc cân bằng tải
	chính xác hơn.
\end{itemize}

Các thông tin MLO được mang trong \textbf{AP MLD TLV} và \textbf{Backhaul STA MLD TLV}
trong bản tin \textit{1905 Topology Response} --- được thiết kế như phần mở rộng tương
thích ngược với các TLV hiện có để hỗ trợ đồng thời các thiết bị Wi-Fi 6 và Wi-Fi 7
trong cùng mạng.


% ============================================================
\section{Duy Trì Danh Sách Tác Nhân (Agent List)}
% ============================================================

Sau khi một Agent mới hoàn tất quá trình Onboarding và cấu hình, Controller có trách
nhiệm đồng bộ thông tin về sự hiện diện của Agent mới này đến \textbf{toàn bộ} các
Agent còn lại trong mạng. Cơ chế này đảm bảo mọi Agent đều có cái nhìn thống nhất về
topo mạng hiện tại --- điều kiện cần thiết để tối ưu hóa lựa chọn đường Backhaul và
phối hợp chuyển vùng liền mạch.

Bản tin \textbf{Agent List Message} mang các thông tin sau cho mỗi Agent trong mạng:
\begin{itemize}
	\item \textbf{AL MAC Address:} Địa chỉ MAC định danh duy nhất của Agent ở lớp
	1905.1 (\textit{AL} --- Abstract Layer MAC).
	\item \textbf{Profile:} Mức Profile bảo mật mà Agent hỗ trợ (Profile-1, -2 hay -3),
	giúp các Agent khác biết cơ chế xác thực nào có thể sử dụng khi thiết lập
	liên kết Backhaul với Agent đó.
	\item \textbf{Backhaul Link Info:} Thông tin về các liên kết Backhaul hiện có của
	Agent (nếu được báo cáo), hỗ trợ các Agent lân cận lựa chọn đường kết nối
	tối ưu.
\end{itemize}

\begin{center}
	\begin{tikzpicture}[>=stealth, font=\small]
		% --- Controller ở trên cùng ---
		\node[draw, fill=gray!15, minimum width=3cm, minimum height=0.65cm]
		(C) at (4.5, 0) {Controller};
		
		% --- Hai Agent bên dưới ---
		\node[draw, fill=blue!8, minimum width=2.8cm, minimum height=0.65cm]
		(A1) at (0, -3.5) {Agent 1 (Mới)};
		\node[draw, fill=gray!10, minimum width=2.8cm, minimum height=0.65cm]
		(A2) at (9, -3.5) {Agent 2 (Hiện có)};
		
		% --- Mũi tên phát quảng bá từ Controller ---
		\draw[->, thick, blue] (C) -- node[scale=0.7, left=2pt, fill=white, inner sep=1pt]
		{Agent List Message} (A1);
		\draw[->, thick, blue] (C) -- node[scale=0.7, right=2pt, fill=white, inner sep=1pt]
		{Agent List Message} (A2);
		
		% --- Nội dung bản tin ---
		\node[draw, dashed, fill=yellow!15, align=left, scale=0.65] at (4.5,-2.0) {
			Nội dung: AL MAC $\cdot$ Profile $\cdot$ Backhaul Link Info
		};
		
		% --- Nhãn giải thích bên dưới mỗi Agent ---
		\node[scale=0.6, gray, align=center] at (0,-4.4)
		{\textit{Nhận danh sách toàn bộ}\\
			\textit{Agent trong mạng}};
		\node[scale=0.6, gray, align=center] at (9,-4.4)
		{\textit{Cập nhật thông tin}\\
			\textit{về Agent mới}};
	\end{tikzpicture}
\end{center}

\textbf{Lưu ý triển khai:} Controller gửi \textit{Agent List Message} đến
\textbf{tất cả} Agent trong mạng --- bao gồm cả Agent mới --- mỗi khi có sự thay đổi
về thành viên (Agent gia nhập hoặc rời mạng). Điều này đảm bảo tính nhất quán topo
tại mọi thời điểm mà không cần cơ chế đồng bộ phân tán.